\scanpage{95}


et on vérifie tout de suite que c'est un hom d'alg. ssi on a

\[ \underbrace{v^2+bv+c}_{\substack{\rotatebox{-90}{$\in$}\\\Gamma L^{\otimes2}}}=0\qquad\text{i.e.}\quad Q(v)=0 \]

On va préciser la relation entre les formes $N$ sur $\mathcal A$, $Q$ sur $E=\mathcal A^\vee$, \deleted{par la forme} \inserted{en utilisant} l'iso (valable pour \unsure{tt|tout} $E$ loc libre de rang 2)

\[ E\xrightarrow{\sim}\underline{\operatorname{Hom}}(E,\underbrace{\det E}_{\substack{\rotatebox{90}{$=$}\mathrm{def}\\\Delta}})\simeq\underset{\substack{\rotatebox{90}{$=$}\mathrm{def}\\\mathcal A}}{E^\vee}\otimes\Delta=\mathcal A\otimes\Delta \]

obtenu en associant à $x\in E$ la forme

\[ f_x:E\unsure{-|\to}\det E=\mathop{\wedge}\nolimits^2 E\qquad y\longrightarrow x\wedge y \]

Si $E$ est muni d'une structure d'extension

\diagram{\begin{tikzcd}[ampersand replacement=\&]0\arrow[r,no head]\&L\arrow[r]\&E\arrow[r,"\psi",no head]\&O\arrow[r]\&0,\end{tikzcd}}

alors

\[ \underset{\substack{\rotatebox{90}{$=$}\\\Delta}}{\det E}\xleftarrow[\sim]{\alpha}L\qquad\text{(avec conv. de signe à préciser)} \]

\[ \alpha(u)=x_0\wedge\unsure{u|v},\quad\text{où }x_0\in\Gamma P_1 \]

\margin{\uline{NB} le $2^{\mathrm e}$ mb ne dépend pas du choix de $x$}

donc

\[ E\xrightarrow{\sim}\underset{\substack{\rotatebox{90}{$=$}\\E^\vee}}{\mathcal A}\otimes L \]

À tout $x\in E$ on associe la forme

\[ g_x:E\longrightarrow L \]

\unsure{donné|donnée} par la condition

\[ g_x(y)=\alpha^{-1}(x\wedge y)\ \deleted{\unsure{x}} \qquad\text{i.e. par la condition} \]

\[ x_0\wedge g_x(y)\deleted{\unsure{x_0|u}}=x\wedge y \]

\deleted{Si $T\in E$ tel que $\psi(T)=2$,} je dis qu'on aura \deleted{pour \unsure{prendre|prouver}}

\[ g_x(y)=\psi(x)y-\psi(y)x \]
% archive sheet number: 94
